图像不是一列没有顺序的特征。相邻像素形成局部结构，同一种边缘可能出现在不同位置。卷积网络把这两个判断写成局部连接和空间参数共享，从而在不显式复制参数的情况下对不同位置使用同一检测规则。

本单元的主线是一次交易的完整账目。写入局部性与平移共享，换来参数量与样本效率；靠深度和空洞把视野推远，换来的视野却在贡献分布上迅速衰减；用下采样买小平移的稳健，赔进去的是空间分辨率和严格的等变性；再用上采样与旁路把分辨率取回来，能取回的只有编码器还留着的那部分。每一步都赔进去了明确的东西，赔的是什么写在算子的结构里，不在超参数里。

\subsection{怎么读这一章}

第一篇是地基：shape 的算术、参数共享和平移等变性，后面四篇都建在它上面。第二篇给出感受野的递推与有效感受野的衰减，是判断``网络能不能看到那个证据''的工具。第三、四篇成对出现，讲分辨率的丢与还，做分割、检测、生成时会直接用到。第五篇不增加新的算子，而是回头给整套先验估价，也是通往注意力单元的接口——如果你正在纠结换不换主干，可以直接从它读起。

读完之后应能区分``参数共享''、``平移等变''和``平移不变''三者，并且知道第三者在真实网络里其实并不成立。

\subsection{本章篇目}

以下篇目按概念依赖排列；若从具体问题进入，也可以先打开最接近当前困惑的一篇，再沿篇内链接回补。

\begin{itemize}
\tightlist
\item \hyperref[sec:14-Deep-Learning/04-Convolutional-Networks/01]{``卷积把平移结构写进线性映射''}。核在所有位置共享同一组权重，卷积因此是一个结构极度规整的 Toeplitz 矩阵；平移等变性是算子自带的代数性质，不需要训练。
\item \hyperref[sec:14-Deep-Learning/04-Convolutional-Networks/02]{``感受野决定了网络能看多远''}。逐层递推算出理论视野，小核堆叠与空洞卷积各自怎样扩大它；理论感受野与有效感受野差一个 \(\sqrt{L}\) 的量级，这是网络``看得到却不看''的原因。
\item \hyperref[sec:14-Deep-Learning/04-Convolutional-Networks/03]{``下采样用分辨率换取不变性''}。池化与步长卷积都是局部聚合接抽稀；不变性只对小于窗口的平移成立，而抽稀在平移量不被步长整除时直接打断等变性，混叠随之而来。
\item \hyperref[sec:14-Deep-Learning/04-Convolutional-Networks/04]{``上采样与分辨率恢复''}。转置卷积是卷积的梯度算子而不是它的逆；核宽不整除步长时覆盖次数不均匀，棋盘纹由此而来；跳跃连接搬运的是解码器在原理上造不出来的信息。
\item \hyperref[sec:14-Deep-Learning/04-Convolutional-Networks/05]{``归纳偏置何时帮忙何时挡路''}。三条先验被做成硬约束，用近似误差换估计误差；CNN 与 ViT 的高下是数据量与偏置强度的交叉点问题，不是先进与否的问题。
\end{itemize}
